\section{Metodología}

\subsection{Búsqueda preliminar}
Se realizó una revisión sistemática de la literatura para vincular la síntesis
de nanopartículas de plata (AgNPs) con sus propiedades ópticas no lineales.
La recolección de datos se ejecutó utilizando bases de datos científicas
principales, incluyendo Scopus, ScienceDirect, Web of Science, SpringerLink y
Google Scholar, y abarcó artículos publicados entre 2020 y 2026.

La búsqueda inicial se centró en artículos de revisión para establecer el
contexto fundamental de la investigación. Los términos generales produjeron un
corpus extenso del cual se extrajo un grupo de trabajos de alta relevancia para
estructurar los antecedentes. Tras establecer este marco, la estrategia se
restringió para aislar artículos de investigación primaria enfocados en la
física fundamental y las propiedades ópticas no lineales de las AgNPs. Se
emplearon las siguientes cadenas de búsqueda booleanas en cada base de datos.

\begin{itemize}
  \item Para Scopus se obtuvieron 249 resultados con la cadena \texttt{TITLE-
      ABS-KEY ((``silver nanoparticle*'' OR ``AgNP*'' OR ``Ag nanoparticle*'') AND
      (``nonlinear optic*'' OR ``nonlinear refractive index'' OR ``nonlinear
    absorption'' OR ``saturable absorption'' OR ``Z-scan'')) AND PUBYEAR AFT 2019}
  \item Para ScienceDirect se obtuvieron 8712 resultados con la cadena
    \texttt{(``silver nanoparticle'' OR AgNP) AND (``nonlinear optics'' OR
    ``nonlinear absorption'' OR ``saturable absorption'' OR z-scan)}
  \item Para SpringerLink se obtuvieron 327 resultados con la cadena
    \texttt{``silver nanoparticles'' AND (``nonlinear optical'' OR ``nonlinear
    absorption'' OR ``Z-scan'')}
  \item Para Google Scholar se obtuvieron 731 resultados con la cadena
    \texttt{intitle:``silver nanoparticles'' OR intitle:``AgNPs'' ``nonlinear
      optics'' OR ``nonlinear optical'' OR ``nonlinear absorption'' OR ``saturable
    absorption'' OR ``Z-scan''}
\end{itemize}

Esta búsqueda dirigida identificó un total combinado de más de \num{10000}
registros preliminares.

\subsection{Limpieza de documentos}
Para refinar el volumen de resultados en un conjunto de datos central, se
aplicó un método de filtración riguroso. El proceso inició con la eliminación
de registros duplicados. Posteriormente, se evaluaron los títulos y resúmenes
para descartar aquellos documentos donde las AgNPs se mencionaban de forma
tangencial o donde los análisis se limitaban a la óptica lineal.

Se integraron investigaciones que presentaran datos experimentales o teóricos
robustos sobre propiedades ópticas no lineales, que investigaran los parámetros
determinantes de dichas propiedades y que se encontraran dentro del periodo de
publicación establecido. Tras esta filtración, el repositorio final se consolidó
en aproximadamente 43 artículos fundamentales que sustentan la discusión crítica
de esta revisión.

\subsection{Análisis bibliométrico}
Para cuantificar el panorama de investigación e identificar tendencias
emergentes, se ejecutó un análisis bibliométrico sobre el conjunto de datos
extraído. Los metadatos bibliográficos pertinentes, tales como autores,
afiliaciones, palabras clave y recuentos de citas, se exportaron desde Scopus
y Web of Science.

Mediante el uso de software de visualización, se generaron mapas de red para
examinar diversas métricas topológicas. Se analizó la coocurrencia de palabras
clave para visualizar el cambio conceptual hacia fenómenos físicos específicos,
como la absorción saturable o el escaneo Z. Asimismo, se evaluaron los patrones
de colaboración institucional y los aumentos repentinos en citaciones para
identificar los trabajos que impulsaron trayectorias de investigación
particulares. Este mapeo validó los grupos temáticos discutidos y evidenció la
transición de la investigación desde la síntesis de nanomateriales hacia
aplicaciones optoelectrónicas avanzadas.
